% generated by roam.py handbook export -- не редактировать, править заметки
% 20260918182821-преобразования_сохраняющие_выпуклость.org (body bdaaf1ab63ff)

\begin{theorem}[{Преобразования, сохраняющие выпуклость}]\label[theorem]{note:6BC7D685-C429-45A3-B563-269840CA0664}
  \begin{enumerate}
    \item \hyperref[note:B9EF495E-CDBE-4E80-B401-8AFEAB0DCCAC]{Аффинные преобразования}: $f: E \rightarrow \R$~--- \hyperref[note:1CA14F2C-637F-41DE-9E97-72075D19E97F]{выпуклая}, $\mathcal{A}(x) = Ax + b$~--- аффинное. Тогда $g(x) = f(Ax + b)$~--- выпуклая на $\left\{ x\ |\ Ax + b \in E \right\}$
    \item $f_{\alpha}$~--- семейство \hyperref[note:1CA14F2C-637F-41DE-9E97-72075D19E97F]{выпуклых функций} на $E$, и $\sup_{\alpha} f_{\alpha}(x) < +\infty\ \forall x \in E$. Тогда $\phi(x) = \sup_{\alpha} f_{\alpha}(x)$~--- выпуклая
    \item $f(x,y): \R^n \times \R^m \rightarrow \R$~--- выпуклая по совокупности переменных $(x, y)$.
      Тогда $\phi(x) = \inf_y f(x,y)$~--- выпуклая (если $\inf$ конечен при всех $x$).
    \item $\Phi: \R^m \rightarrow \R$~--- выпуклая и неубывающая по каждому аргументу, $f_1, \dots, f_m$~--- выпуклые.
      Тогда $\Phi(f_1(x), \dots, f_m(x))$~--- выпуклая
  \end{enumerate}
\end{theorem}

\begin{proof}
  Везде $\lambda \in [0;1]$, $z = \lambda x + (1-\lambda) y$, проверяется определение \hyperref[note:1CA14F2C-637F-41DE-9E97-72075D19E97F]{выпуклой функции}.

  \proofpart{1. Аффинное преобразование}
  Пусть $Ax + b, Ay + b \in E$. \hyperref[note:B9EF495E-CDBE-4E80-B401-8AFEAB0DCCAC]{Аффинное отображение} сохраняет выпуклые комбинации: $Az + b = \lambda (Ax + b) + (1-\lambda)(Ay + b)$, так что $Az + b \in E$ (область определения $g$ выпукла) и
  \[
  g(z) = f(\lambda (Ax + b) + (1-\lambda)(Ay + b)) \le \lambda f(Ax + b) + (1-\lambda) f(Ay + b) = \lambda g(x) + (1-\lambda) g(y)
  \]

  \proofpart{2. Супремум семейства}
  Для каждого $\alpha$: $f_{\alpha}(z) \le \lambda f_{\alpha}(x) + (1-\lambda) f_{\alpha}(y) \le \lambda \phi(x) + (1-\lambda) \phi(y)$, так как $f_{\alpha} \le \phi$ поточечно.
  Правая часть не зависит от $\alpha$, поэтому можно взять $\sup$ по $\alpha$ слева: $\phi(z) \le \lambda \phi(x) + (1-\lambda) \phi(y)$.

  \proofpart{3. Инфимум по части переменных}
  Возьмем $\varepsilon > 0$ и $y_1, y_2$ такие, что $f(x, y_1) < \phi(x) + \varepsilon$, $f(y, y_2) < \phi(y) + \varepsilon$ (существуют по определению $\inf$).
  Точка $(z,\ \lambda y_1 + (1-\lambda) y_2) = \lambda (x, y_1) + (1-\lambda)(y, y_2)$, и по выпуклости $f$ по совокупности переменных
  \[
  \phi(z) \le f(z,\ \lambda y_1 + (1-\lambda) y_2) \le \lambda f(x, y_1) + (1-\lambda) f(y, y_2) < \lambda \phi(x) + (1-\lambda) \phi(y) + \varepsilon
  \]
  Устремим $\varepsilon \rightarrow 0$.

  \proofpart{4. Композиция с монотонной выпуклой функцией}
  Обозначим $F(x) = (f_1(x), \dots, f_m(x))$. По выпуклости каждой $f_k$: $F(z) \le \lambda F(x) + (1-\lambda) F(y)$ покоординатно.
  $\Phi$ не убывает по каждому аргументу, поэтому $\Phi(F(z)) \le \Phi(\lambda F(x) + (1-\lambda) F(y))$, а по выпуклости $\Phi$ правая часть $\le \lambda \Phi(F(x)) + (1-\lambda) \Phi(F(y))$.
\end{proof}
